Con el fin de alcanzar los objetivos descritos anteriormente, es necesario definir el alcance del sistema y los elementos organizativos que se verán afectados por su implantación.

El sistema desarrollado en este proyecto está orientado principalmente a apoyar la actividad comercial de un distribuidor de productos capilares profesionales y a mejorar la interacción con sus clientes, que en este caso son peluquerías profesionales.

La aplicación se plantea como una plataforma web accesible desde distintos dispositivos, diseñada con un enfoque responsive y compatible con su uso en dispositivos móviles mediante tecnología PWA (Progressive Web Application). De esta forma, tanto el distribuidor como las peluquerías podrán acceder al sistema desde ordenadores, tablets o teléfonos móviles.

Desde el punto de vista organizativo, el sistema afecta principalmente a tres tipos de usuarios:

\begin{itemize}
    \item El distribuidor, responsable de la gestión de pedidos, la organización de rutas de reparto, el seguimiento de entregas y la gestión administrativa asociada a su actividad comercial.
    \item Los comerciales que operan bajo la responsabilidad del distribuidor, quienes podrán utilizar el sistema para gestionar su actividad diaria, registrar visitas a clientes, realizar pedidos en nombre de las peluquerías y consultar información relevante para el desarrollo de su actividad comercial.
    \item Las peluquerías clientes, que podrán consultar el catálogo de productos, realizar pedidos y acceder a distintas herramientas digitales relacionadas con la gestión de su actividad profesional.
\end{itemize}

El alcance del sistema se limita al ámbito de la relación comercial entre el distribuidor y las peluquerías profesionales. Por tanto, el sistema no contempla en esta fase inicial funcionalidades orientadas a clientes finales o consumidores particulares, aunque se deja abierta la posibilidad de ampliar el sistema en el futuro para incorporar este tipo de servicios.